\documentclass[10pt,a4paper]{amsart}
\usepackage[margin=19mm]{geometry}
\usepackage{amsmath,amssymb,mathtools}
\usepackage[scr=rsfs,cal=euler]{mathalfa}
\usepackage{xfrac}
\usepackage[unicode,hidelinks,bookmarks=false]{hyperref}
\newcommand{\ie}{i.e.\ }
\newcommand{\cf}{cf.\ }
\newcommand{\resp}{respectively\ }
\newcommand{\Subsection}{\subsection}
\providecommand{\nobreakdash}{\nobreak\hbox{-}}
\setlength{\emergencystretch}{2em}
\multlinegap0pt
\usepackage[shortlabels]{enumitem}
\usepackage{mathtools}
\usepackage{amssymb}
\newtheorem{thm}{Theorem}
\newtheorem{proposition}{Proposition}
\DeclareMathOperator{\val}{val}
\title{Equidistribution for Semiabelian Varieties over Number Fields}
\author{Wei Xue}
\date{Source: arXiv:2608.08262v2 (CC BY 4.0)}
\begin{document}
\maketitle

\noindent\emph{Source and licence.} Equidistribution for Semiabelian Varieties over Number Fields. Version arXiv:2608.08262v2 (CC BY 4.0), distributed by arXiv under the Creative Commons Attribution 4.0 International licence, http://creativecommons.org/licenses/by/4.0/, as recorded in the arXiv licence metadata for this exact version and shown on the abstract page https://arxiv.org/abs/2608.08262v2. Full licence notice: \texttt{licenses/arxiv-2608.08262-CC-BY-4.0.txt}.\par\smallskip

\noindent\textit{Editorial scope.} Counted unit: the article's thm thm:local-trivialization (source0.tex lines 1452--1486). The article's complete original proof of it is delivered with it (source0.tex lines 1892--2017, one \texttt{proof} environment of 126 lines, which the article places away from the statement). No mathematics is written, repaired or re-derived: the statement and the proof are the article's own lines, in the article's own notation, and no retained line is substituted or re-flowed.  Retained uncounted as the source-local context the proof needs: lines 1868--1890.  Nothing was omitted from any retained span, so the package carries no omission record.  No retained line contains a citation, so the wrapper carries no bibliography and no bibliography entry.  No retained line contains a figure environment, an image-inclusion command or drawing code, so no image is delivered and the package declares no asset dependency.  Editorial formatting only, in addition: an AMS \texttt{amsart} preamble that restates the article's own declarations and macros used by the retained lines, namely the environments \texttt{thm}, \texttt{proposition} on separate counters, together with the standard packages those lines need.  The retained environments print in this document as Theorem 1, Proposition 1 in order of appearance, whereas the article prints the same items with the numbering of its own full document.  The complete native source of the article as distributed in the arXiv e-print for this exact version is bundled unchanged as \texttt{sources/207144/source0.tex}.
\begin{proposition}[Analytic inputs for the local trivialization]\label{input:local-triv-analytic}
	The proof of Theorem~\ref{thm:local-trivialization} uses the following standard
	facts in the precise forms stated here.
	\begin{enumerate}[leftmargin=2em]
		\item Raynaud--Bosch--Lutkebohmert uniformization gives
		\[
		0\to M'_v\to E_v^{an}\xrightarrow{p_v}A_{\mathbb C_v}^{an}\to0
		\]
		and the associated Raynaud extension.  Over a sufficiently small contractible
		analytic open \(U\subset A_{\mathbb C_v}^{an}\), the covering map \(p_v\) admits
		an analytic section.
		\item For a rational section \(s\) of a rigidified algebraically trivial line
		bundle on \(A\), Raynaud theta descent produces a function \(f_{s,v}\) satisfying
		\[
		\operatorname{div}(f_{s,v})=p_v^{-1}(\operatorname{div}(s)^{an}).
		\]
		\item A morphism from the principal torus to a normal toric variety extends over
		the affine chart \(X_\sigma\) precisely when all characters
		\(\chi^m\), \(m\in\sigma^\vee\cap M\), pull back to analytic functions with
		effective boundary divisor.  This is the standard affine-chart criterion for
		toric morphisms, applied in the analytic category.
	\end{enumerate}
\end{proposition}

\begin{thm}[Non-archimedean local trivialization]\label{thm:local-trivialization}
	Let \(v\) be a non-archimedean place of \(K\).  There is a finite collection
	\(\{(U_j,\psi_j)\}_{j\in J}\) such that the \(U_j\) form an open covering of
	\(A^{an}_{\mathbb C_v}\), and
	\[
	\psi_j:
	\overline{G}^{an}_{\mathbb C_v}|_{U_j}
	\longrightarrow
	X_{\Sigma,\mathbb C_v}^{an}
	\]
	are analytic maps for which
	\[
	\overline{\pi}^{an}_{\mathbb C_v}\times\psi_j:
	\overline{G}^{an}_{\mathbb C_v}|_{U_j}
	\longrightarrow
	U_j\times X_{\Sigma,\mathbb C_v}^{an}
	\]
	is an isomorphism of analytic spaces.  Moreover, for every toric
	\(\mathbb R\)-Cartier divisor \(D\),
	\[
	\overline{G}(D)^{\mathrm{can}}_v|_{U_j}
	=
	\psi_j^*(\overline{D}^{\mathrm{can}}_v).
	\]
	For all \(j,j'\in J\) and all
	\[
	y\in \overline{G}^{an}_{\mathbb C_v}|_{U_j\cap U_{j'}},
	\]
	one has
	\[
	\operatorname{val}_v\circ\psi_j(y)
	=
	\operatorname{val}_v\circ\psi_{j'}(y).
	\]
\end{thm}

\begin{proof}[Proof of Theorem~\ref{thm:local-trivialization}]
	Let \(\Sigma(1)=\{\rho_1,\ldots,\rho_d\}\), let \(v_\ell\) be the primitive
	generator of \(\rho_\ell\), and let \(D_\ell\) be the corresponding
	\(T\)-invariant prime divisor on \(X_\Sigma\).  Write the extension class of
	\(G\) in the chosen character basis as
	\[
	\eta_G=(H_1,\ldots,H_t)\in A^\vee(\overline K)^t .
	\]
	The cocycle description of the associated toric fibration gives
	\[
	\overline{\pi}^*H_i
	=
	\sum_{\ell=1}^{d}\langle e_i^\vee,v_\ell\rangle\,\overline G(D_\ell),
	\qquad 1\le i\le t.
	\]
	
	Choose an analytic open subset \(U\subset A_{\mathbb C_v}^{an}\) which is
	contractible and disjoint from the analytic divisors of rational sections
	\(s_i\) of \(H_i^\vee\).  By Proposition~\ref{input:local-triv-analytic}, the
	covering \(p_v:E_v^{an}\to A_{\mathbb C_v}^{an}\) has an analytic section over
	\(U\); denote the corresponding lift of \(\pi\) by
	\[
	\widetilde\pi_U:\overline G^{an}_{\mathbb C_v}|_U\longrightarrow E_v^{an}.
	\]
	The section \(s_i\) determines a theta function \(f_{s_i,v}\) with
	\[
	\operatorname{div}(f_{s_i,v})
	=
	p_v^{-1}\bigl(\operatorname{div}(s_i)^{an}_{\mathbb C_v}\bigr).
	\]
	The cocycle formula gives a rational function \(g_i\) on \(\overline G\) with
	\[
	\operatorname{div}(g_i)
	=
	\sum_{\ell=1}^{d}\langle e_i^\vee,v_\ell\rangle\,\overline G(D_\ell)
	+\overline\pi^*\operatorname{div}(s_i).
	\]
	After multiplying \(s_i\) by constants, we may assume
	\(g_i(e_G)=1\) and \(\Vert f_{s_i,v}(e_E^{an})\Vert=1\).
	
	On the open semiabelian part over \(U\), define
	\[
	\psi_U^{(i)}(y)
	=
	\frac{g_i(y)}{f_{s_i,v}(\widetilde\pi_U(y))},
	\qquad
	\psi_U=(\psi_U^{(1)},\ldots,\psi_U^{(t)}):G^{an}_{\mathbb C_v}|_U\to T^{an}_{\mathbb C_v}.
	\]
	The denominator has no zero over \(U\), and hence \(\psi_U^{(i)}\) is an
	invertible analytic function on \(G^{an}_{\mathbb C_v}|_U\).
	For each \(a\in U\), the divisor of \(\psi_U^{(i)}\) on the compactified
	fiber \(\overline G_a\simeq X_{\Sigma,\mathbb C_v}\) is exactly the divisor of
	the torus character \(\chi^{e_i^\vee}\).  Thus \(\psi_U\) identifies the
	principal torus in the fiber with \(T_{\mathbb C_v}^{an}\), up to multiplication
	by nonzero constants.
	
	We now extend over the toric boundary.  Let
	\(X_\sigma=\operatorname{Spec}\mathbb C_v[\sigma^\vee\cap M]\) be an affine
	toric chart.  For \(m=\sum_i m_ie_i^\vee\in\sigma^\vee\cap M\), the pullback
	of \(\chi^m\) is
	\[
	\prod_{i=1}^{t}\bigl(\psi_U^{(i)}\bigr)^{m_i}.
	\]
	Its boundary divisor is
	\[
	\sum_{\ell=1}^{d}\langle m,v_\ell\rangle\,\overline G(D_\ell),
	\]
	which is effective on the chart because \(m\in\sigma^\vee\cap M\).  By the
	affine-chart criterion in Proposition~\ref{input:local-triv-analytic}, the map
	\((\overline\pi,\psi_U)\) extends uniquely to an analytic morphism
	\[
	\Psi_U:
	\overline G^{an}_{\mathbb C_v}|_U
	\longrightarrow
	U\times X_{\Sigma,\mathbb C_v}^{an}.
	\]
	The same character calculation applied to the inverse torus coordinates gives
	an inverse morphism.  Hence \(\Psi_U\) is an isomorphism of analytic spaces.
	
	The finite cover is obtained by choosing finitely many such contractible opens
	and, on each of them, rational sections \(s_i\) whose divisors avoid the open.
	Compactness of \(A_{\mathbb C_v}^{an}\) gives a finite subcover
	\(\{U_j\}_{j\in J}\), and we write the second component of \(\Psi_{U_j}\) as
	\(\psi_j\).
	
	It remains to verify overlap compatibility.  On \(U_j\cap U_{j'}\), write
	\[
	s_j^{(i)}=t_{jj'}^{(i)}s_{j'}^{(i)}
	\]
	with \(t_{jj'}^{(i)}\) an analytic unit.  The corresponding rational functions
	and theta functions satisfy
	\[
	g_j^{(i)}=(t_{jj'}^{(i)}\circ\pi)g_{j'}^{(i)},\qquad
	f_{s_j^{(i)},v}=(t_{jj'}^{(i)}\circ p_v)f_{s_{j'}^{(i)},v}.
	\]
	If the chosen local lifts of \(U_j\) and \(U_{j'}\) differ on a connected
	component of the overlap by an element of the period lattice \(M'_v\), the
	theta transformation law adds the usual Raynaud automorphy factor.  The same
	factor occurs in the descent cocycle defining the semiabelian extension
	class \(\eta_G\); hence it cancels in the quotient
	\(g_i/f_{s_i,v}\).  After this cancellation the displayed formulas give the
	actual change of local coordinates.
	Since \(p_v(\widetilde\pi_U(y))=\pi(y)\), the two unit factors have the same
	absolute value after evaluation at \(y\).  Therefore
	\[
	\left|
	\frac{\psi_j^{(i)}(y)}{\psi_{j'}^{(i)}(y)}
	\right|_v=1,\qquad 1\le i\le t.
	\]
	Equivalently, the transition multiplier lies in the compact torus, so it has
	valuation zero.  The toric tropicalization map is invariant under compact
	torus multiplication.  Hence
	\[
	\val_v\circ\psi_j(y)=\val_v\circ\psi_{j'}(y)
	\]
	on overlaps.
	
	Finally, under \(\Psi_U\), the line bundle \(\overline G(D)\) is identified
	with the pullback of the toric line bundle \(O(D)\) on the second factor.  When
	\(D\) is equipped with the canonical toric metric, that metric is a function
	of the toric valuation and is invariant under compact torus multiplication.
	The overlap compatibility therefore identifies the canonical metrics:
	\[
	\overline G(D)^{can}_v|_{U_j}=\psi_j^*(\overline D^{can}_v).
	\]
\end{proof}

\end{document}
